\section{The four-way input audit}
\label{sec:input-audit}

\paragraph{Scope and meaning of the deliverable.}
\texttt{INPUT\_CLASSIFICATION.tsv} contains one primary class for each
of the 362 rows of the supplied GRH\_NIGHT
\texttt{CONSEQUENCES.tsv}, 162 ALPOGE\_BRIDGES findings, and all 39
CX\_GRH findings described by the attached intake.  The ALPOGE set
includes the intake's full numeric ranges and contextual shorthand,
together with an explicitly marked conservative supporting block.
Question/current-gate rows are recorded separately in
\texttt{INPUT\_UNCLASSIFIED\_GATES.tsv}; a missing question statement is
not assigned a fictitious class \emph{none}.  The supplied CX intake
reports earlier proof checks, but the original CX proof package was not
available or resolved in this audit.  Those 39 classifications are
therefore explicitly \emph{intake-only}.

This is a dependency and provenance audit.  A row's original statement,
input, status, proof locator and Grothendieck locator are preserved
verbatim for the GRH catalogue.  Targeted source readings resolve the
ambiguous or additional inputs discussed below.  Other inherited proof
claims are carried as background with that exact check status; the table
does not certify a new complete reading of 563 mathematical proofs.

\paragraph{Conventions.}
The four values are exactly \texttt{A}, \texttt{M},
\texttt{outside preprint}, and \texttt{none}.  A denotes Alp\"oge's
Dirichlet-GRH theorem, used as a theorem as instructed.  M denotes the
conditional claim for every unitary Hecke character over every number
field in Masullo, \S1.2--\S1.2.3, pp.~12--16.  Its finite-order part
suffices for Brauer/Artin and Dedekind applications; potentially CM
positive-weight factors generally require infinite-order unitary Hecke
characters after the weight twist.  Theorem 13.3.1.2, p.~522, is not used
as a theorem statement.

The single-column precedence is
\[
 \texttt{M}>\texttt{outside preprint}>\texttt{A}>\texttt{none}.
\]
It is an administrative convention, not an ordering of mathematical
implication or necessity.  In particular, M does not prove PF26's
algebraic theorem.  The outside-input names and all additional premises
remain in separate columns.  A smaller exact analytic family is retained
even when M is its sufficient source.  The class \texttt{none} means no
use of A, M, or an identified substantive outside-preprint theorem in the
assertion being classified.  It does not mean no literature, no
hypotheses, or an unconditional solution of every surrounding problem.
For example NT-L15 and NT-Z19 retain the stronger weak mean-square
hypothesis, and NT-Z07 retains its independent non-CM $\mathrm{GL}_2$
GRH hypothesis even though its primary class is A.

\subsection{MA-IN-001: the three misleading leading A labels}

\paragraph{Finding.}
The full asserted scope of GRH\_NIGHT NT-Z13, NT-Z15 and NT-Z16 has
primary class M.  Their Dirichlet specializations have class A.  The
source catalogue's opening A label is weaker than the scope of the
complete row; the remainder of each original input field already
acknowledges the conditional Hecke branch.

\begin{proof}
NT-Z13 applies its argument both to zeta/Dirichlet functions and to
finite-order Hecke characters over general number fields.  Its stated
analytic hypotheses include the real or conjugate-paired gamma
parameters; the arbitrary infinite-order case is explicitly excluded.
NT-Z15 applies the central-multiplicity bound to the same Hecke scope.
NT-Z16 invokes the NT-Z13 analytic hypotheses for its first
\emph{nonnegative} ordinate bound.  In all three, the zero-location
input for the general Hecke branch is M's finite-order assertion.
A's exported Dirichlet assertion does not itself supply that branch.
Thus the union of actually asserted cases has primary class M under the
declared rule.  Direct inspection of the complete source input fields
shows that M is mentioned after the opening A text, so the correction
is to a misleading leading label rather than an allegation that the
source completely omitted M.  NT-Z14 inherits the same distinction and
is also classified M, without the misleading-leading-label flag.
\end{proof}

In geometric language, a finite-order Hecke character gives a rank-one
local system over its own number field.  Induction to $\mathbb Q$ need
not split into rank-one Dirichlet systems.  This is the Artin-motive
scope distinction used in GRH\_NIGHT AG-ARTIN and MT-ART01, and in
ALPOGE\_BRIDGES CC-F03; the finite-\'etale representation dictionary is
SGA~1, Expos\'e~V, \S7.  It is not a logical independence assertion
between the two GRH statements.

\subsection{MA-IN-002: the PF26 input in AG-CYCLIC-DIV}

\paragraph{Finding.}
GRH\_NIGHT AG-CYCLIC-DIV has class \texttt{outside preprint}, with
PF26 named.  Its original word ``neither'' correctly excludes A and M
but is insufficient as a four-way input label.  The full original input
field already names PF26; no undisclosed-dependency allegation is made.

\begin{proof}
The supplied geometry proof invokes PF26's odd-splitter construction in
\S4: Lemma~4.1, p.~15; Lemma~4.2, pp.~15--16; Lemmas~4.3--4.4,
p.~17; Corollary~4.5, p.~18; and Proposition~4.6, pp.~18--19.
These are affirmative algebraic inputs to the proof, not papers merely
named as targets of criticism.  No A/M zero-location input is needed for
this supplied algebraic step.  Consequently the required class is
\texttt{outside preprint}.  The separate full PF26 reading in MA-PF-001
checks this source chain; this documentary argument neither duplicates
that proof nor turns PF26 into a formally certified theorem.
\end{proof}

The geometric content is computation of a nontrivial idempotent in a
finite \'etale algebra, hence a decomposition of a finite \'etale
scheme (SGA~1, Expos\'e~V, \S7).  GRH\_NIGHT CX-ALG-027 and
CX-ALG-028 need both the algebraic PF26 input and an analytic family.
They have primary class M at full generality, with PF26 and the precise
A/pure-field premises explicitly retained.  When the bad-prime family
of MA-PF-002 is empty, the latter specialization uses PF26+A; the table
does not conceal this specialization under universal M.

\subsection{MA-IN-003: abelian motives and the named 2026 inputs}

\paragraph{Finding.}
The cp03 ALPOGE intake's ``none''/``unconditional'' description of the
full abelian-sector conclusions CC-AG01--03 and CC-XC01--04 must retain
the outside-preprint input van Hoften--Sempliner,
Hexagon~2610.00259v2.  The later intake's explicit preprint qualification
is the correct provenance.  CC-TATE03 has a second distinct input,
Hexagon~2610.00226v1.

\begin{proof}
The supplied module \texttt{consequences\_criteria\_Tate2026.tex}
states the chain explicitly.  CC-TATE01 imports van Hoften and
Sempliner's \emph{The Hodge Conjecture for CM abelian varieties},
Theorem~I, p.~1, version~2 of 18 October 2026.  CC-TATE02 applies
Milne's published 1999 Theorem~7.1, journal p.~72, and Propositions
7.3--7.4, pp.~73--75, together with his published 2002 Theorem~3.3
and Corollary~3.4, final author-manuscript pp.~8--9.  Those classical
implications assume precisely the all-CM Hodge input.  The abelian
numerical-motive classifications and arbitrary-abelian forward-cycle
targets in the flagged rows use this conclusion.  Neither A nor M
supplies the algebraic cycle theorem.

CC-TATE03 separately imports \emph{Milne's rationality conjecture},
Hexagon~2610.00226v1, contributed by Keerthi Madapusi and bearing the
PDF byline ``ChatGPT and Claude'': Theorem~1.1, p.~1;
Theorem~5.1, pp.~11--12; Proposition~A.1, pp.~13--14; and
Corollary~6.2, p.~13.  The rationality assertion uses this manuscript;
the simultaneous algebraic cycle on the special fibre additionally
uses the all-CM theorem.  Both are retained.  Milne's
\emph{On the Hodge, Tate, and standard conjectures for abelian
varieties}, Hexagon~2611.00244v1, Theorems~2.1, 3.1 and 3.5,
pp.~4, 6 and 8, is recorded as the cited 2026 corroborating note.
It is not counted as an indispensable third theorem where the proof
already cites the original 1999/2002 implications.  This establishes
the precise dependency correction.
\end{proof}

The relevant Grothendieck question is the passage from Frobenius
eigenvalues to actual motivic objects and correspondences in cote~15,
scans~100, 102, 104 and 106, as located in ALPOGE\_BRIDGES
CC-TATE04.  A spectral coincidence is not itself an algebraic
correspondence.  The table also preserves the important narrower
exceptions: CC-XC20--22 use their explicit Tate/supersingular Chow
decompositions and published Fu--Li results without the all-abelian
preprint.  BR-15-CLASSBOUND and BR-15-SLOPEBOUND use classical
Honda--Tate and ideal theory.  They are not assigned the broader input
merely because they occur in the same programme.

\subsection{MA-IN-004: a single class preserves a conjunction only with its notes}

\paragraph{Finding.}
The table's single primary label is compatible with all required
conjunctions and conditional subcases precisely because the secondary
inputs and assertion scope are retained.  A formal equivalence and an
application asserting its positive or zero-location side receive
different descriptions.

\begin{proof}
For a row let $S$ be its set of identified A/M/outside-theorem inputs
and $H$ its remaining hypotheses.  The primary label is the maximum of
$S$ under the declared administrative precedence, or \texttt{none}
if $S$ is empty.  The stored original input, outside-name list and
residual notes retain $S$ and $H$; thus the projection is never used to
replace the conjunction by its maximum.  In particular the projection
does not imply $\mathrm{M}\Rightarrow\mathrm{PF26}$.

For the formal criteria MT-M19, MT-M20 and MT-M27, the proof establishes
an equivalence without asserting that its GRH/positivity side holds.
These rows have class \texttt{none}; the A/M application is stated
separately, and simplicity remains additional in MT-M20.  In contrast,
AG-DET, AG-MOMENT, AG-HANKEL and AG-SIGNED explicitly include the
identified arithmetic Hecke/CM determinant or quotient application.
Their positive real squared-zero support uses M at full scope, or A
in the Dirichlet specialization.  Their elementary entire-function,
moment or Gram-matrix cores can be used with independently supplied
positive support, but that observation does not erase the arithmetic
input of the full row.  Likewise MT-M26 retains M for its asserted
every-positive-width narrow-strip application, whereas its wide-strip
core is unconditional.  These examples establish that the convention
does not convert an arithmetic zero-location premise into ``none''
merely by calling it a hypothesis.
\end{proof}

For cote~1, scans~79--80 and 111--112, the distinction is between the
general Fredholm/trace formalism and positivity of the specific
arithmetic operator $B$.  The source framework alone does not supply
its positive spectrum.  ALPOGE\_BRIDGES BR-TR-01--02 and
GRH\_NIGHT MT-M20/M23 preserve the further distinction between positive
traces and faithful positive full-jet operators.

\subsection{MA-IN-005: the coverage certificate}

\paragraph{Finding.}
There are 563 distinct classified run/ID pairs.  All 362 supplied
GRH catalogue rows occur exactly once, with their five original source
fields unchanged.  Every primary class is one of the required four
values.  The independent range/context audit passes.

\begin{proof}
The deterministic builder uses explicit reviewed ID assignments, not
an A/M keyword rule.  It refuses duplicate assignments or a mismatch
between the source catalogue and the assignment set.  The independent
script \texttt{checks/audit\_input\_classification.py} reads the emitted
TSV and original source files without importing the builder's decisions.
It verifies equality of the 362-ID sets, uniqueness of run/ID pairs,
exact preservation of statement/input/status/proof/Grothendieck fields,
the permitted class values, nonempty evidence/status/residual fields,
existing evidence files and named outside inputs.  It separately
expands the intake's numeric ranges (126 finding IDs), checks 22
contextual nonnumeric IDs, verifies the 39-entry CX finding set, and
requires all six CX gate IDs in the separate gate register.
The emitted certificate \texttt{checks/input\_independent\_audit.json}
records PASS and the table's SHA-256.  This proves the finite data
claims just stated; it does not prove the inherited mathematics.
\end{proof}

\begin{center}
\begin{tabular}{lrrrrr}
Run & A & M & outside preprint & none & total\\
\hline
GRH\_NIGHT &66&140&3&153&362\\
ALPOGE\_BRIDGES &25&0&50&87&162\\
CX\_GRH &8&14&5&12&39\\
\hline
Total &99&154&58&252&563
\end{tabular}
\end{center}

\subsection{MA-IN-006: other preprints, published sources and refutation targets}

\paragraph{Finding.}
The outside-preprint class must retain substantive imported manuscripts
beyond the four names listed in the prompt.  Publication status and
logical use both matter: an arXiv URL alone does not determine the class,
and refuting a manuscript does not import its theorem as a premise.

\begin{proof}
Direct dependency readings identify Engel--de Gaay Fortman--Schreieder,
arXiv:2507.15704v3, Theorem~1.1 and Corollary~1.2, p.~2, in AG-DP-EGS
and the very-general even-degree branch of AG-FOURIER-ISOGENY.
Their institution lists the work as a preprint.  CC-XC23 imports
Fu--Huang--Li, arXiv:2609.38801v2, Theorem~1.9, p.~5, with its
excellent-lifting, characteristic and Artin-invariant restrictions.
CC-WM01 imports Ito--Zou, arXiv:2610.07890v1, Theorem~B, p.~2,
and Wear, arXiv:2303.05610v1, Theorem~1.1.2, p.~3, in two of its
branches.  Wear's author page still lists the paper as submitted.
CC-WM02--03 import Petrov--Zavyalov,
arXiv:2607.14517v2, Theorem~1.2, p.~2, equivalently
Theorem~5.1 and Lemma~5.2, p.~18.  F-SM-32--33 use
Ramzi--Sosnilo--Winges, arXiv:2503.11338v1,
Definition~2.5 and Theorem~2.6(1),(4), pp.~5--6; the author lists
acceptance in JEMS, but an issued publication was not established in
this audit.  The table retains it as an accepted preprint with that
qualification.  CX-E01--05 use the separate OpenAI result E,
\emph{The Exponential Time Hypothesis is False}, dated 25 October 2026,
Theorem~1.1 and Corollary~1.4 as reported by the intake.  The official
repository locates it among preprints.  These are substantive premises
and are therefore named.

Conversely, primary publisher records establish that Deninger's
arXiv:2504.15767 result is published in \emph{Journal of Number Theory}
286 (2026), 47--66; Haine--Porta--Teyssier's exodromy paper is published
in \emph{Journal of Topology} 19(3) (2026); and the partial-regularity
theorem used by F-SM-15 is published as Chen, Lau, Lee, Ono and Zhang,
\emph{Archiv der Mathematik} 127 (2026), 165--178.
These are retained as published inputs.  Finally BKNO
arXiv:2508.17776, Corollary~3.2 and Proposition~4.17, is the target
of the explicit counterexample in AG-LOCAL-EXCEPTION; Thakur's
Proposition~2.9 is the target in BR-15-CLASSBOUND.  The counterexamples
do not assume the assertions they refute.  This establishes all the
claimed distinctions; exact URLs and dated evidence appear in
\texttt{modules/inputs\_source\_status.tsv}.
\end{proof}

\subsection{MA-IN-007: the exact Alp\"oge locator and its inspection boundary}

\paragraph{Finding.}
The supplied runs support the following precise citation for A:
Alp\"oge, \emph{Remark}, public four-page announcement,
\url{https://alpo.ge/remark.pdf}, p.~4, step~20; the reported interface
is \texttt{ALPOGE\_Main.lean}, with definitions in
\texttt{ALPOGE\_Statements.lean}.  This audit uses A as a theorem and
makes no fresh Lean-compilation claim.

\begin{proof}
The public PDF's p.~4 identifies step~20 as GRH.  The predecessor's
\path{research/source_and_masullo.md}, \S1.1, records its complete
reading of the following exported declarations:
\begin{verbatim}
riemannHypothesis : RiemannHypothesis := rhOfAll
riemannZeta_eq_zero_re : RiemannZetaZerosOnLine := rhOfAll
LFunction_ne_zero_of_half_lt_re
  : DirichletLNonvanishingRightOfHalf := allChi
LFunction_eq_zero_re : DirichletLZerosOnLine := claimAll.2.1
LFunction_eq_zero_re_of_isPrimitive
  : PrimitiveDirichletLZerosOnLine := critAll
\end{verbatim}
The same record identifies the pin
\texttt{0df444a360eaa60ab8c11dca51a86af692955474}
and the import \texttt{Supplements.Kernels.MainTheoremsAll}.
Its definitions establish the rational Dirichlet scope; no exported
arbitrary-number-field Hecke assertion is recorded there.
The present attachment contains the predecessor's inspection record,
not those Lean files or their full import tree.  Thus the exact filenames
and declarations can be cited as predecessor-inspected interface
evidence, while the public step locator was checked directly.  No new
formal-certificate event follows from either fact.
\end{proof}
